\section{Counting-model transfer and a reachability-summary distinction}

\subsection{A simultaneous normalization for every deletion subset}
\label{app:counting}

The forest counting theorem of Enright, Meeks, and
Molter~\citep{enright2025counting} uses undirected temporal paths with
nondecreasing timestamps. The following explicit transformation explains
why fixed-deletion counting transfers to the present directed, strict,
parallel-event convention. The transformation below establishes the comparison;
it is not asserted as a theorem stated in the cited work.

\begin{proposition}[Deletion-preserving counting normalization]
For an explicit depth-three temporal in-arborescence with $M>0$ events, there
is a polynomial-time transformation to a temporal forest with at most $M$
distinct time labels such that, simultaneously for every retained event
subset, the number of temporal paths from vertices of original depth $h$
to the root is exactly $W_h$, for $h=2,3$. Original event identities are
preserved bijectively.
\end{proposition}

\begin{proof}
Let $\rho(e)$ be the rank of $t(e)$ among distinct original timestamps,
preserving ties. Let $d(e)\in\{1,2,3\}$ be the original depth of the source
of $e$, and give each original event a distinct index
$i(e)\in\{0,\ldots,M-1\}$. Define
\[
 t'(e)=4M\rho(e)+Md(e)+i(e).
\]
For consecutive rootward events $e,f$, their depths differ by one, so
\[
 t'(e)-t'(f)=4M(\rho(e)-\rho(f))+M+i(e)-i(f).
\]
The final term lies between $1$ and $2M-1$. Hence
$t(e)<t(f)$ if and only if $t'(e)\le t'(f)$. Equal original timestamps
are incompatible in the rootward direction. All transformed labels are
distinct: different ranks are separated by the $4M$ term, different depths
within a rank by $M$, and different identities within a rank and depth by
their indices. Rank-compressing these labels preserves their order and
leaves at most $M$ time labels.

Forget arc directions. The underlying graph is a tree, so the unique simple
path from any vertex to the root follows the original rootward direction.
A temporal path cannot leave that route and return without repeating a
vertex. Its successive transformed labels are nondecreasing exactly when
the original labels are strictly increasing. Distinct original parallel
contacts remain distinct time-edges. The correspondence survives deletion
of any common subset of identities and preserves the number of events in
each path. Summing over starting vertices of original depth $h$ gives the
claimed equality.
\end{proof}

Empty histories are handled directly. Forced statuses and costs may be
carried as identity metadata, but the cited counting theorem does not
optimize them. Counting for each fixed deletion set therefore transfers;
enumerating and optimizing a shared deletion set still requires a separate
argument. This normalization concerns simple temporal paths on a forest,
not arbitrary undirected walks that may return along an edge.

\subsection{Reached sets and earliest arrivals do not encode tuple quotas}
\label{app:reachability}

Consider the path $a\to v\to u\to z$, with three optional first events at
times $1,4,7$, three retained middle events at times $2,5,8$, and a retained
final event at time $10$. There are three two-walks and six three-walks.
Deleting the first event at time one or the first event at time four leaves,
respectively, three or four triples; both choices leave all three pairs and
have the same unit deletion cost.

Reverse every arc and replace time $t$ by $11-t$, taking $z$ as the source
and assigning traversal time one. For either deletion, the same four vertices
are reachable with earliest-arrival times $(0,2,4,5)$ along the reversed
path. The final reversed contact at time four comes from the undeleted
original first event at time seven in both cases. Yet the original upper
quota $W_3\le3$ accepts only the first deletion. No statistic consisting
solely of these reached sets and earliest arrivals can distinguish the two
choices, even after assigning fixed vertex weights.

This example separates the indicated summary from tuple multiplicities.
It does not prove that reachability methods cannot be extended, nor that an
auxiliary encoding is impossible.
